\chapter{Matrices, Vectorized Computing, and Linear Algebra with NumPy}
\label{chap:matrices-numpy}

\section{Why NumPy: Overcoming Python Loop Latency}
Python lists store pointers to individual \texttt{PyObject} wrappers scattered across the heap, incurring continuous cache misses and dynamic type-checking overhead. NumPy introduces the contiguous C-level \texttt{ndarray}, unlocking CPU SIMD vector registers and BLAS acceleration.

\begin{figure}[htbp]
\centering
\begin{tcolorbox}[colback=boxgraybg,colframe=boxnavyborder,title={Memory Comparison: List of Floats vs NumPy ndarray}]
\small
\textbf{Python List (3 floats):} \texttt{[Ptr0, Ptr1, Ptr2]} $\longrightarrow$ 3 separate 24-byte \texttt{PyFloatObject}s scattered on heap.\\
\textbf{NumPy ndarray (3 floats):} Contiguous 24-byte buffer of IEEE 754 raw doubles in contiguous physical memory.
\end{tcolorbox}
\caption{Cache Locality: Python Dynamic Pointers vs Contiguous C Buffer}
\end{figure}

\section{\texttt{ndarray} Internal Architecture and Striding}
\begin{definition}[Array Strides]
The \textbf{strides} of an array are a tuple of integers specifying the number of bytes to jump in memory to advance by one index along each axis.
\begin{equation}
\text{Element Pointer}(i_0, i_1) = \text{Data Buffer Base} + i_0 \cdot s_0 + i_1 \cdot s_1
\end{equation}
For a C-contiguous array of shape $(M, N)$ storing 64-bit floats ($8$ bytes per element):
\begin{equation}
s_0 = N \cdot 8 \text{ bytes}, \qquad s_1 = 8 \text{ bytes}
\end{equation}
Transposing a matrix simply swaps the strides tuple: $\mathcal{O}(1)$ time with zero memory copying!
\end{definition}

\section{Broadcasting Rules and Formal Algebra}
\begin{theorem}[NumPy Broadcasting Theorem]
Two dimensions are compatible when:
\begin{enumerate}
    \item They are equal in size, OR
    \item One of them is 1.
\end{enumerate}
If an axis has size 1, NumPy virtualizes expansion along that axis by setting its stride to 0!
\end{theorem}

\section{Fundamental Matrix Operations and Inverses}
\begin{lstlisting}
import numpy as np

# Create matrices
A = np.array([[3.0, 1.0], [1.0, 2.0]])
b = np.array([9.0, 8.0])

# Direct linear system solve: Ax = b
x = np.linalg.solve(A, b)
print(f"Solution vector x: {x}")  # Output: [2., 3.]

# Matrix product (@ operator)
residual = A @ x - b
print(f"Norm of residual: {np.linalg.norm(residual):.2e}")
\end{lstlisting}

\section{Matrix Factorization in Data Science}
\begin{table}[htbp]
\centering
\small
\begin{tabularx}{\linewidth}{lp{3.5cm}X}
\toprule
\textbf{Decomposition} & \textbf{Matrix Requirements} & \textbf{Machine Learning \& Engineering Application} \\
\midrule
\textbf{Cholesky} ($A = L L^T$) & Symmetric Positive-Definite & Efficient Gaussian process sampling, Kalman filtering ($\frac{1}{3} n^3$ FLOPs). \\
\textbf{QR} ($A = Q R$) & General Rectangular & Numerically stable ordinary least squares regression without squaring condition number. \\
\textbf{SVD} ($A = U \Sigma V^T$) & Any Real Matrix $M \times N$ & Principal Component Analysis (PCA), low-rank approximation, latent semantic analysis. \\
\textbf{Eigen} ($A V = V \Lambda$) & Square Matrix & Spectral clustering, graph PageRank, Markov chain stationary distributions. \\
\bottomrule
\end{tabularx}
\caption{Core Matrix Factorization Algorithms in SciPy / NumPy}
\end{table}

\section{Chapter Summary}
\begin{summarybox}
$\bullet$ Zero-copy operations (slices, transpositions) manipulate stride metadata rather than allocating memory.\\
$\bullet$ Avoid Python-level loops; leverage vectorized ufuncs to achieve near-C compute speeds.\\
$\bullet$ Prefer \texttt{np.linalg.solve(A, b)} over explicit matrix inversion \texttt{np.linalg.inv(A) @ b}.
\end{summarybox}
